\exnumber{03}
\extitle{UX II - Ex 4 - Problem statement}
\exfiles{ex03/PDF, named Name\_FirstName\_DeliverableName\_Date\_VersionNumber}
\exauthorize{Those named in the exercise}
\makeheaderfiles

\textbf{Instructions:}

Based on your previous work and your empathy map, define a problem statement that would clarify the main problem you want to tackle.

You may use the How Might We exercise in pairs.

\textbf{Deliverable:}

\begin{itemize}
\item A PDF document with a problem statement starting with a "How Might We"
\item Upload the PDF document on Github, labeled as follows: "Name\_FirstName\_DeliverableName\_Date\_VersionNumber". For example, it can be labeled as "Doe\_John\_MidFiPrototype\_05242024\_V1".
\item Add a title within the document.
\item Pay attention to the global aspect of the document. It must be clear, simple and easy to read. You can get inspiration on canva.com but don't overload your design with too much details!
\end{itemize}

\textbf{Resources:}

\begin{itemize}
\item Stuck on a problem? Just ask how might we \url{https://relab.academy/design-thinking/stuck-on-a-problem-just-ask-how-might-we/}
\item Problem statement \url{https://www.indeed.com/career-advice/career-development/what-is-a-problem-statement}
\item Using 'How Might We' Questions to Ideative on the Right Problems \url{https://www.nngroup.com/articles/how-might-we-questions/}
\end{itemize}

\textbf{Quote:}

'Out of clutter, find simplicity. From discord, find harmony. In the middle of difficulty, lies opportunity.' Albert Einstein

